% GENERATED FILE. Edit canonical lessons, then run npm run generate:books.
% canonical-source-hash: fnv1a64:3326886974b42d36
% canonical-lessons: GU-C298-mrtyu, GU-C298-sok, GU-C298-aghat, GU-C298-asvasan, GU-C298-sahevu

\chapter{A Sad Day}
\label{ch:gu-a-sad-day}
\hlchaptermodality{298}

\begin{chapteropening}
\textbf{By the end of this chapter:} \emph{I can speak about a death and comfort a family.}

Complete the last lesson of chapter 298: I can speak about a death and comfort a family.
\end{chapteropening}

\section[\texorpdfstring{\gu{મૃત્યુ} (mṛtyu)}{મૃત્યુ (mṛtyu)}]{\gu{મૃત્યુ} (mṛtyu) — death}

\label{lesson:GU-C298-mrtyu}



\begin{quote}
Before the new one: say the Gujarati for a gathered crowd, then the Gujarati for a push.
\end{quote}

\subsection*{You'll want to know: \gu{મૃત્યુ}}
\textbf{\gu{મૃત્યુ}} — \emph{mṛtyu} — \textquotedblleft{}death\textquotedblright{}.

\begin{quote}
\gu{અકસ્માતમાં} \gu{એક} \gu{માણસનું} \gu{મૃત્યુ} \gu{થયું}.
\end{quote}

\begin{quote}
\emph{akasmātmā̃ ek māṇasnũ mṛtyu thayũ.}
\end{quote}

\begin{quote}
One man died in the accident.
\end{quote}

A neuter noun.

\begin{cousinweb}
\textbf{\gu{મૃત્યુ}} (\emph{mṛtyu}) is a tatsama, Sanskrit \emph{mṛtyu}, \textquotedblleft{}death, dying\textquotedblright{} (Monier-Williams). It is the word of news reports: \textbf{\gu{તેનું} \gu{મૃત્યુ} \gu{થયું}}, he died.
\end{cousinweb}

\begin{grammarlens}[title={What you've built}]
One more word for telling what happened.
\end{grammarlens}

\subsection*{Your turn}
\begin{itemize}
\raggedright
  \item \emph{Say it:} \emph{mṛtyu}
  \item \emph{Say it:} \emph{mṛtyu}, once more
  \item \emph{Say it:} \emph{akasmātmā̃ ek māṇasnũ mṛtyu thayũ.} — One man died in the accident.
  \item \emph{Recall:} say the Gujarati for a gathered crowd, then the Gujarati for a push, then say \emph{mṛtyu} again
\end{itemize}

\begin{tcolorbox}[breakable,colback=teal!4,colframe=teal!35!black,title={Before you move on}]
What does \gu{મૃત્યુ} mean? (\textquotedblleft{}Death\textquotedblright{}.) Say it once more.
\end{tcolorbox}
\section[\texorpdfstring{\gu{શોક} (śok)}{શોક (śok)}]{\gu{શોક} (śok) — mourning; grief}

\label{lesson:GU-C298-sok}



\begin{quote}
Before the new one: say the Gujarati for a push, then the Gujarati for death.
\end{quote}

\subsection*{You'll want to know: \gu{શોક}}
\textbf{\gu{શોક}} — \emph{śok} — \textquotedblleft{}mourning; grief\textquotedblright{}.

\begin{quote}
\gu{આખા} \gu{ગામમાં} \gu{શોક} \gu{છે}.
\end{quote}

\begin{quote}
\emph{ākhā gāmmā̃ śok che.}
\end{quote}

\begin{quote}
The whole village is in mourning.
\end{quote}

A masculine noun.

\begin{cousinweb}
\textbf{\gu{શોક}} (\emph{śok}) is a tatsama, Sanskrit \emph{śoka}, \textquotedblleft{}burning, sorrow\textquotedblright{} (Monier-Williams). It is grief that is shared and shown: \textbf{\gu{શોક} \gu{પાળવો}}, to keep mourning.
\end{cousinweb}

\begin{grammarlens}[title={What you've built}]
One more word for telling what happened.
\end{grammarlens}

\subsection*{Your turn}
\begin{itemize}
\raggedright
  \item \emph{Say it:} \emph{śok}
  \item \emph{Say it:} \emph{śok}, once more
  \item \emph{Say it:} \emph{ākhā gāmmā̃ śok che.} — The whole village is in mourning.
  \item \emph{Recall:} say the Gujarati for a push, then the Gujarati for death, then say \emph{śok} again
\end{itemize}

\begin{tcolorbox}[breakable,colback=teal!4,colframe=teal!35!black,title={Before you move on}]
What does \gu{શોક} mean? (\textquotedblleft{}Mourning; grief\textquotedblright{}.) Say it once more.
\end{tcolorbox}
\section[\texorpdfstring{\gu{આઘાત} (āghāt)}{આઘાત (āghāt)}]{\gu{આઘાત} (āghāt) — a shock (of bad news)}

\label{lesson:GU-C298-aghat}



\begin{quote}
Before the new one: say the Gujarati for death, then the Gujarati for mourning.
\end{quote}

\subsection*{You'll want to know: \gu{આઘાત}}
\textbf{\gu{આઘાત}} — \emph{āghāt} — \textquotedblleft{}a shock (of bad news)\textquotedblright{}.

\begin{quote}
\gu{દીકરાના} \gu{મૃત્યુથી} \gu{માતાને} \gu{મોટો} \gu{આઘાત} \gu{લાગ્યો}.
\end{quote}

\begin{quote}
\emph{dīkrānā mṛtyuthī mātāne moṭo āghāt lāgyo.}
\end{quote}

\begin{quote}
Her son's death came as a terrible shock to the mother.
\end{quote}

A masculine noun.

\begin{cousinweb}
\textbf{\gu{આઘાત}} (\emph{āghāt}) is a tatsama, Sanskrit \emph{āghāta}, \textquotedblleft{}striking, a blow\textquotedblright{} (Monier-Williams): bad news that hits like a blow. It \textbf{\gu{લાગે} \gu{છે}}, strikes you: \textbf{\gu{મોટો} \gu{આઘાત} \gu{લાગ્યો}}.
\end{cousinweb}

\begin{grammarlens}[title={What you've built}]
One more word for telling what happened.
\end{grammarlens}

\subsection*{Your turn}
\begin{itemize}
\raggedright
  \item \emph{Say it:} \emph{āghāt}
  \item \emph{Say it:} \emph{āghāt}, once more
  \item \emph{Say it:} \emph{dīkrānā mṛtyuthī mātāne moṭo āghāt lāgyo.} — Her son's death came as a terrible shock to the mother.
  \item \emph{Recall:} say the Gujarati for death, then the Gujarati for mourning, then say \emph{āghāt} again
\end{itemize}

\begin{tcolorbox}[breakable,colback=teal!4,colframe=teal!35!black,title={Before you move on}]
What does \gu{આઘાત} mean? (\textquotedblleft{}A shock (of bad news)\textquotedblright{}.) Say it once more.
\end{tcolorbox}
\section[\texorpdfstring{\gu{આશ્વાસન} (āśvāsan)}{આશ્વાસન (āśvāsan)}]{\gu{આશ્વાસન} (āśvāsan) — comfort; consolation}

\label{lesson:GU-C298-asvasan}



\begin{quote}
Before the new one: say the Gujarati for mourning, then the Gujarati for a shock.
\end{quote}

\subsection*{You'll want to know: \gu{આશ્વાસન}}
\textbf{\gu{આશ્વાસન}} — \emph{āśvāsan} — \textquotedblleft{}comfort; consolation\textquotedblright{}.

\begin{quote}
\gu{પડોશીઓએ} \gu{કુટુંબને} \gu{આશ્વાસન} \gu{આપ્યું}.
\end{quote}

\begin{quote}
\emph{paḍośīoe kuṭumbne āśvāsan āpyũ.}
\end{quote}

\begin{quote}
The neighbours comforted the family.
\end{quote}

A neuter noun.

\begin{cousinweb}
\textbf{\gu{આશ્વાસન}} (\emph{āśvāsan}) is a tatsama, Sanskrit \emph{āśvāsana}, \textquotedblleft{}reviving, consoling\textquotedblright{} (Monier-Williams). It is given: \textbf{\gu{આશ્વાસન} \gu{આપ્યું}}, offered comfort.
\end{cousinweb}

\begin{grammarlens}[title={What you've built}]
One more word for telling what happened.
\end{grammarlens}

\subsection*{Your turn}
\begin{itemize}
\raggedright
  \item \emph{Say it:} \emph{āśvāsan}
  \item \emph{Say it:} \emph{āśvāsan}, once more
  \item \emph{Say it:} \emph{paḍośīoe kuṭumbne āśvāsan āpyũ.} — The neighbours comforted the family.
  \item \emph{Recall:} say the Gujarati for mourning, then the Gujarati for a shock, then say \emph{āśvāsan} again
\end{itemize}

\begin{tcolorbox}[breakable,colback=teal!4,colframe=teal!35!black,title={Before you move on}]
What does \gu{આશ્વાસન} mean? (\textquotedblleft{}Comfort; consolation\textquotedblright{}.) Say it once more.
\end{tcolorbox}
\section[\texorpdfstring{\gu{સહેવું} (sahevũ)}{સહેવું (sahevũ)}]{\gu{સહેવું} (sahevũ) — to bear; to endure}

\label{lesson:GU-C298-sahevu}



\begin{quote}
Before the new one: say the Gujarati for a shock, then the Gujarati for consolation.
\end{quote}

\subsection*{You'll want to know: \gu{સહેવું}}
\textbf{\gu{સહેવું}} — \emph{sahevũ} — \textquotedblleft{}to bear; to endure\textquotedblright{}.

\begin{quote}
\gu{માતાએ} \gu{મોટું} \gu{દુઃખ} \gu{સહ્યું}.
\end{quote}

\begin{quote}
\emph{mātāe moṭũ duḥkh sahyũ.}
\end{quote}

\begin{quote}
The mother bore a great sorrow.
\end{quote}

Past, with \gu{તેણે} (he or she …): \gu{તેણે} \gu{દુઃખ} \gu{સહ્યું}, \emph{teṇe duḥkh sahyũ}.

\begin{cousinweb}
\textbf{\gu{સહેવું}} (\emph{sahevũ}) is inherited from Sanskrit \emph{sáhatē}, \textquotedblleft{}conquers; bears\textquotedblright{}, and CDIAL lists the Gujarati verb, \textquotedblleft{}bears\textquotedblright{} (CDIAL 13304). Its past is \textbf{\gu{સહ્યું}}, as \textbf{\gu{કહેવું}} gives \textbf{\gu{કહ્યું}}.
\end{cousinweb}

\begin{grammarlens}[title={What you've built}]
One more word for telling what happened.
\end{grammarlens}

\subsection*{Your turn}
\begin{itemize}
\raggedright
  \item \emph{Say it:} \emph{sahevũ}
  \item \emph{Say it:} \emph{sahevũ}, once more
  \item \emph{Say it:} \emph{mātāe moṭũ duḥkh sahyũ.} — The mother bore a great sorrow.
  \item \emph{Recall:} say the Gujarati for a shock, then the Gujarati for consolation, then say \emph{sahevũ} again
\end{itemize}

\begin{tcolorbox}[breakable,colback=teal!4,colframe=teal!35!black,title={Before you move on}]
What does \gu{સહેવું} mean? (\textquotedblleft{}To bear; to endure\textquotedblright{}.) Say it once more.
\end{tcolorbox}
